% 导读：用平实的话说明这门课在讲什么、各讲之间的关系、几种收敛的区别
\frontsection{导读：这门课在讲什么}

这一部分不讲证明，只说清楚三件事：这门课最终要得到什么结论；为了得到这些结论，中间要经过哪些站；每一站上出现的名词是什么意思。读正文时如果迷失了方向，回到这里看一看自己在哪一站。

\subsection*{一、一个贯穿全书的例子：反复抛硬币}

把一枚公平的硬币抛 $n$ 次，记正面的次数为 $S_n$。大家在中学就知道，$n$ 很大时，正面的频率 $S_n/n$ 应该接近 $\frac12$。这门课的核心内容，就是把这句话说准确，并且说得更多。

\textbf{(1) 大数定律：频率趋于 $\frac12$。}可是“趋于”是什么意思？$S_n/n$ 是随机的，抛出一万次正面也不是不可能，只是概率极小。所以“趋于”有两种合理的理解：
\begin{itemize}
\item \emph{弱的说法}：对固定的很大的 $n$，$S_n/n$ 与 $\frac12$ 相差超过 $0.01$ 的\emph{概率}很小，并且 $n\to\infty$ 时这个概率趋于零。这叫\term{依概率收敛}，这个结论叫\term{弱大数定律}。
\item \emph{强的说法}：想象你真的一直抛下去，得到一个无穷的结果序列 $\omega$（比如“正反反正正……”），对这一个具体的序列计算频率，它作为一个数列收敛到 $\frac12$。除了概率为零的“坏序列”（比如一直是正面），每个序列都如此。这叫\term{几乎必然收敛}，这个结论叫\term{强大数定律}。
\end{itemize}

\textbf{(2) 中心极限定理：偏差的形状是正态曲线。}大数定律说 $S_n/n-\frac12\to0$，但没有说偏差有多大、什么形状。计算表明 $S_n-\frac n2$ 的大小约为 $\sqrt n$ 的量级，于是看标准化之后的
\[
Z_n=\frac{S_n-n/2}{\sqrt n/2}.
\]
中心极限定理说：$n$ 很大时，$Z_n$ 的直方图越来越像标准正态分布的钟形曲线，即
\[
P(Z_n\leqslant x)\to\Phi(x)=\int_{-\infty}^x\frac1{\sqrt{2\pi}}e^{-t^2/2}\dd t .
\]
请注意：这里\emph{不是} $Z_n$ 本身收敛到某个随机变量。对一个具体的抛硬币序列，$Z_1,Z_2,Z_3,\dots$ 会一直上下乱跳，永远不会稳定下来。收敛的只是它的\emph{分布}，也就是“$Z_n$ 落在各个区间里的概率”。这就是\term{依分布收敛}（也叫\term{弱收敛}），记作 $Z_n\Rightarrow N(0,1)$。

\textbf{(3) 再往后：鞅与 Markov 链。}如果每次抛硬币的规则依赖于过去（比如赌徒根据输赢调整赌注，或者明天的天气依赖于今天的天气），各次结果就不再独立。鞅和 Markov 链是处理这种“有记忆的随机过程”的两个基本模型。课程的最后三分之一讲它们。

所以整门课的主线是：
\[
\begin{aligned}
&\text{概率空间}\ \to\ \text{随机变量与分布}\ \to\ \text{期望}\ \to\ \text{独立性}\\
&\quad\to\ \text{大数定律}\ \to\ \text{中心极限定理}\ \to\ \text{条件期望与鞅}\ \to\ \text{Markov 链}.
\end{aligned}
\]

\subsection*{二、为什么要先学测度论}

初等概率论里，“抛 $n$ 次硬币”只有 $2^n$ 种结果，每种的概率是 $2^{-n}$，一切都是有限的求和。可是强大数定律说的是“一直抛下去”，样本空间是全体无穷 $0$-$1$ 序列，它是不可数集，每个具体序列的概率都是 $0$。这时“频率收敛到 $\frac12$ 的那些序列的概率是 $1$”这句话是什么意思？要回答它，必须先回答：
\begin{itemize}
\item 哪些集合可以谈论概率？（不是所有集合都行，这引出 \term{$\sigma$-代数}，第~\ref{lec:set-systems} 讲。）
\item 概率要满足什么规则？（可数可加性，这引出\term{测度}，第~\ref{lec:measures} 讲。）
\item 满足这些规则的概率真的存在吗？比如 $[0,1]$ 上的“长度”。（这是\term{测度的延拓}，第~\ref{lec:extension-of-measures} 讲。）
\item 期望 $EX$ 在无穷的、连续的情形下怎样定义？（这是 \term{Lebesgue 积分}，第~\ref{lec:integration} 讲。）
\end{itemize}
测度论就是为了回答这些问题而建立的。它本身不是概率论的目的，而是概率论的语言。很多测度论的名词，在概率论里只是换了个名字。下面这张对照表请同学们熟记：

\begin{center}
\small
\begin{tabular}{p{14.5em}p{10.5em}p{11em}}
\toprule
测度论的说法 & 概率论的说法 & 意思\\
\midrule
测度空间 $(\Omega,\cF,\mu)$，$\mu(\Omega)=1$ & 概率空间 $(\Omega,\cF,P)$ & 全部可能的结果，及其概率\\
$\cF$ 中的集合 & 事件 & 可以问“发生了没有”的事情\\
可测函数 & 随机变量 & 由试验结果决定的数\\
像测度 $\mu\circ X^{-1}$ & 分布 $\mu_X$ & $X$ 落在各处的概率\\
积分 $\int X\dd\mu$ & 期望 $EX$ & 加权平均\\
几乎处处（a.e.） & 几乎必然（a.s.） & 除了一个概率为零的事件之外\\
乘积测度 & 独立 & 互不影响的试验\\
Radon--Nikodym 导数 & 密度；以后的条件期望 & 一个测度相对另一个的“比率”\\
\bottomrule
\end{tabular}
\end{center}

\subsection*{三、几种收敛到底有什么区别}

这是初学者最容易糊涂的地方，我们用最简单的例子把它们分开。随机变量 $X_n$ 是 $\Omega$ 上的函数。数学分析里，函数列有逐点收敛、一致收敛、积分意义下的收敛；概率论里也有好几种，此外还多了一种只看分布的收敛。

\begin{enumerate}
\item \textbf{几乎必然收敛} $X_n\to X$ a.s.：对几乎每一个样本点 $\omega$，数列 $X_n(\omega)$ 收敛到 $X(\omega)$。这就是“几乎处处的逐点收敛”。\\
\emph{例子}：强大数定律，$S_n/n\to\frac12$ a.s.
\item \textbf{依概率收敛} $X_n\to X$ in probability：对每个 $\varepsilon>0$，$P(\abs{X_n-X}>\varepsilon)\to0$。也就是说，对每个大的 $n$，“$X_n$ 与 $X$ 相差较大”这件事的概率很小。它比几乎必然收敛弱：允许每个 $\omega$ 都“偶尔”偏离一下，只要偏离的时刻越来越稀少。\\
\emph{例子}：弱大数定律。
\item \textbf{$L^p$ 收敛} $X_n\to X$ in $L^p$：$E\abs{X_n-X}^p\to0$。这是“平均意义下”的收敛。\\
\emph{例子}：$E(S_n/n-\frac12)^2=\frac1{4n}\to0$。
\item \textbf{依分布收敛} $X_n\Rightarrow X$：$P(X_n\leqslant x)\to P(X\leqslant x)$（在 $X$ 的分布函数的连续点 $x$ 处）。它只比较\emph{分布}，根本不比较 $X_n$ 与 $X$ 的值；$X_n$ 与 $X$ 甚至可以定义在不同的概率空间上。\\
\emph{例子}：中心极限定理。
\end{enumerate}

它们之间的关系是
\[
\begin{array}{ccccc}
\text{几乎必然收敛}&\Longrightarrow&\text{依概率收敛}&\Longrightarrow&\text{依分布收敛}\\
&&\Uparrow&&\\
&&L^p\ \text{收敛}&&
\end{array}
\]
反方向的箭头一般都不成立，第~\ref{lec:convergence} 讲会给出每一个反例，并且告诉读者这些反例为什么要那样构造。这里先给一个最能说明“依分布收敛很弱”的例子：设 $X$ 以概率 $\frac12$ 取 $1$，以概率 $\frac12$ 取 $-1$，令 $X_n=-X$ 对一切 $n$。则每个 $X_n$ 与 $X$ 同分布，所以 $X_n\Rightarrow X$；可是 $\abs{X_n-X}=2$ 处处成立，$X_n$ 离 $X$ 一点也不近。依分布收敛只说“$X_n$ 的直方图越来越像 $X$ 的直方图”，至于 $X_n$ 与 $X$ 的值是否接近，它什么也不管。

\subsection*{四、讲次路线图}

每一讲都在回答一个问题。下表列出各讲回答的问题。

\begin{center}
\small
\begin{tabular}{p{1.2em}p{12.3em}p{20.5em}}
\toprule
讲 & 主题 & 回答的问题\\
\midrule
\ref{lec:set-systems} & 集合系，$\pi$-$\lambda$ 定理 & 哪些集合可以谈论概率？怎样证明“一族集合都具有某个性质”？\\
\ref{lec:measures} & 测度及其唯一性 & 概率满足什么规则？什么时候两个概率测度相同？\\
\ref{lec:extension-of-measures} & 测度的延拓 & 只在区间上给出的“长度”能否扩充成全体 Borel 集上的测度？\\
\ref{lec:random-variables} & 随机变量 & 什么是随机变量？哪些运算之后仍是随机变量？\\
\ref{lec:distributions} & 分布与分布函数 & 怎样描述一个随机变量“取各个值的概率”？\\
\ref{lec:integration} & 积分 & 怎样对一般的函数求平均？极限与积分何时可以交换？\\
\ref{lec:absolute-continuity} & 绝对连续与密度 & 一个分布什么时候有密度？\\
\ref{lec:expectation} & 期望 & 怎样计算期望？期望与概率之间有哪些不等式？\\
\ref{lec:product-measures} & 乘积测度，Fubini 定理 & 怎样描述两个互不影响的试验？二重积分何时可以交换次序？\\
\ref{lec:independence} & 独立性 & 独立性的精确含义是什么？怎样判断？怎样构造无穷多个独立随机变量？\\
\ref{lec:convergence} & 几种收敛方式 & 随机变量列“收敛”有哪些不同的意思？它们之间什么关系？\\
\bottomrule
\end{tabular}
\end{center}
其后依次是：弱大数定律，Borel--Cantelli 引理与 $0$-$1$ 律，强大数定律，弱收敛的进一步理论，特征函数，中心极限定理，Poisson 收敛；条件期望，鞅，鞅收敛定理，可选停时定理；Markov 链，常返与暂留，平稳分布。这些讲写好之后会补进这张表。

\subsection*{五、怎样读这份笔记}

\begin{enumerate}
\item 每一讲先读开头的引言与末尾的“本讲小结”，知道这一讲要去哪里，再读中间。
\item 第一遍只读定义、定理的陈述、例子与例题，跳过标有 ★ 的部分以及较长的证明，但要读证明前的“证明思路”。第二遍再读证明。
\item 每遇到一个新定义，先在例子中把它具体算一遍：比如遇到依概率收敛，就拿抛硬币的 $S_n/n$ 去验证。抽象的定义只有在具体的例子中才能理解。
\item 读到一个定理，问自己：每个条件去掉会怎样？正文中通常有对应的反例。
\item 作业的解答紧跟在题目后面。可以先想五分钟，想不出就读解答，读完合上，自己再写一遍。
\end{enumerate}
